\paragraph{The height of the boundary columns.}
The arithmetic proof of Proposition~\ref{prop:geometric-run-index}
must use one bound for each of the two determinant costs
in \eqref{eq:geometric-four-terms}. The boundary determinant
below is evaluated in the trace dual. The complementary
determinant is evaluated by CM coefficient functionals.
\hypertarget{source-slice-budget}{}
\begin{lemma}[Short trace-dual tests for boundary columns]
Let \(r_z>0\), \(d=R_j\) and \(n=[K:\mathbb Q]\). For every
\(B>d\sqrt n\exp(P_LS_\kappa+P_d/2)\) the \(K\)-dual of the
column space has a spanning family of functionals \(\phi\) satisfying the two bounds
\[
 \phi(M)\subseteq\mathfrak D_K^{-1},\qquad |\sigma\phi(v)|\le B\|v\|_{2,\sigma}\quad(\text{every }\sigma).
\]

\end{lemma}
\begin{proof}
Equip \(M\) with the Euclidean norm from the sum of the
conjugate-compatible squared \(L^2\) norms. Its real dimension
is \(nd\). The height bound \(H(v)\ge-P_LS_\kappa\) for integral
\(v\ne0\), the inequalities
\(\|v\|_{2,\sigma}\le\|v\|_{\sup,\sigma}
 \le e^{P_d/2}\|v\|_{2,\sigma}\), and arithmetic--geometric mean give
\[
 \lambda_1(M)\ge\sqrt n\,e^{-P_LS_\kappa-P_d/2},\qquad
 \lambda_{nd}(M^*)\le \frac{nd}{\lambda_1(M)}<B .
\]
The second inequality is Banaszczyk's transference theorem.
By nondegeneracy of trace the rational functional attached to
an integral dual vector is \(v\mapsto{\rm Tr}_{K/\mathbb Q}\phi(v)\).
Testing on \(av\) for every \(a\in\mathcal O_K\) gives the
codifferent condition. Each Hermitian component of its Riesz
vector has norm at most its full Euclidean norm. Applying these
facts to the spanning short dual vectors gives both assertions.
\cite{banaszczyk}\label{cor:slice-budget-1}
\end{proof}

\begin{proposition}[The boundary determinant cost]\label{prop:slice-boundary-determinant-cost}
For the selected \(r_z\) boundary columns and the norm bounds
above, the finite costs \(G_\lambda\) have finite support and satisfy
\[
 \langle G\rangle\le r_z H_z S_\kappa,\qquad H_z=2G_c+P_{rw}+P_L+\tfrac32(P_d+P_O).
\]

If there are no columns, the cost is zero.
\end{proposition}

\begin{proof}
\proofstep{The determinant of the boundary columns.}
The boundary columns are independent over \(F\), so some
minor of their values on these functionals is a nonzero
polynomial. Choose one nonzero coefficient \(A\in K\).
At each prime, multiplying all its rows by an element of
order \(v_\lambda(\mathfrak D_K)\) makes the functionals
integral there. The definition of \(G_\lambda\) gives
\[
 0\le G_\lambda\le v_\lambda(A)+r_z v_\lambda(\mathfrak D_K).
\]
This proves finite support and supplies one coefficient for
all primes, though the minimizing coefficient may vary.

For each embedding, expand that coefficient multilinearly
in the coefficient vectors \(z_{i,\gamma}\) of the columns.
Every determinant is a pairing of exterior products, bounded
by \(B^{r_z}\prod_i\|z_{i,\gamma_i}\|_{2,\sigma}\)
by Cauchy--Schwarz and Hadamard. Summing bounds over
\emph{all} tuples of indices therefore gives
\[
 \log |\sigma A|\le r_z\log B+
       \sum_i\log\left(\sum_\gamma\|z_{i,\gamma}\|_{2,\sigma}\right).
\]
For column \(i\) write \(a=T-s\). Its coefficient sum is
at most \((N e^{m_\sigma})^a\|t_i^0\|_{\sup,\sigma}\),
by expanding the light products, where \(m_\sigma\) bounds
all the light log suprema. Its mean is bounded in logarithm
by \(a(\log N+G_c S_L)+G_c S_0\). The weight size convention
gives \[
 \begin{gathered}
 aS_L+S_0=S_\kappa+a, \\ a\le S_\kappa, \\ a\log N\le P_{rw}a(S_L-1)\le P_{rw}S_\kappa.
\end{gathered}
\]
Use the product formula for \(A\) and
\(\mathrm N\mathfrak D_K=|\operatorname{disc}K|\).

Let \(B\) decrease to the transference bound. Since
\(\log d\le P_d\), \(\log n\le P_O\) and
\(\log|\operatorname{disc}K|/n\le P_O\), this proves
the asserted \(\langle G\rangle\le r_z H_z S_\kappa\).
For no boundary columns the determinant is one and the price
is zero.

\end{proof}

\paragraph{The complementary quotient and its CM determinant.}

The boundary-column estimate leaves the determinant of a complementary quotient. Its torsion-free lattice and CM evaluation give the second determinant contribution.

Write \(Q=M/(M\cap H_g)\) and \(\rho=r_q\). It is torsion-free
over \(\mathcal O_K\). Choose \(\theta_1,\ldots,\theta_\rho\)
out of the short spanning sections of weight \(\kappa\) to be
a rational basis modulo \(H_g\). Here the short-section price
for each mean log norm is at most \(H_\theta\). We explain
the coindex of their images. At the definite end \(\rho=0\)
and the price is zero.

\begin{proposition}[The complementary determinant cost]\label{prop:slice-complement-determinant-cost}
For the chosen \(r_q\) basis sections of the complementary
quotient, assume the CM coefficient and line bounds below.
For the stated \(P_Q\), the weighted quotient length satisfies
\[
 \langle L_Q\rangle\le2r_qP_Q S_\kappa.
\]
 The coefficient argument uses the same value-line
coordinates at finite and archimedean places. At the definite
lift the quotient and its cost are zero.
\end{proposition}

\begin{proof}
\proofstep{Coefficient functionals on the complementary quotient.}
Before the last lift, the functionals used in the coefficient
test extract a monomial in formal frame letters and a Laurent
coefficient on a CM arc. Evaluation is linear in the section.
The frame polynomial is homogeneous of degree \(A\), the sum of the nonnegative flag exponents, so a nonzero test has that total monomial degree. Its value is
measured in a one-dimensional space with the dual theta
generator to power \(n_\kappa\) (the central weight) and the dual differential
generators to the powers in this monomial. Thus an arc
functional \(\phi_i\) has offsets \(s_{iw}\) at finite places
and logarithmic norm offsets \(a_{i\sigma}\) at infinity.
The finite offsets have finite support. Over its coefficient
field the three CM estimates used are
\begin{align*}
 w(\phi_i(v))+s_{iw}&\ge-e_{w/\lambda} b_\lambda,\\
 \log|\sigma\phi_i(v)|+a_{i\sigma}
    &\le \log\|v\|_{\sup,\sigma|K}+C_{\rm arch},\\
 \bigg(\sum_w s_{iw}\log {\rm N}w
       -\sum_{\sigma} a_{i\sigma}\bigg)/[K':\mathbb Q]
    &\le B_{\rm deg}.
\end{align*}

In this formula \(v\in M\) has nonzero reading, and \(\sigma\)
runs over embeddings. The last inequality concerns the
degree of the value line. In particular its two terms must
use the same differential and theta coordinates. These
bounds apply to any conjugate of the test and persist
after a number-field extension: finite orders and their offsets
multiply by relative ramification, and the archimedean
data pull back along embeddings. The common baseline for
the allowed torsion order \(n'\) is
\[
 b_\lambda=A\max\{v_\lambda(|\operatorname{disc}E|),
                               v_\lambda(n')\}.
\]
It is nonnegative and its mean is at most
\(A(\log|\operatorname{disc}E|+\log n')\) by norming the
two principal ideals. A change of generator of the value
line changes its two offsets by a principal divisor, so
the degree inequality is unaffected.

The CM annihilation and detection of low coefficients give
common rational kernel \(H_g\) in the rational cusp space. To select independent rows on its quotient, include
all coefficient-field conjugates of the tests. Over \(K\)
the coordinate readings of a test in any basis of its finite
value field are linear combinations over a normal closure
of the conjugate readings, by the invertible embedding
matrix.
For a value field of degree \(n\), a basis \(b_1,\ldots,b_n\) and its distinct embeddings \(\sigma_i\), the inversion uses
\[
  \det\bigl(\sigma_i(b_j)\bigr)_{i,j=1}^n\ne0 .
\]
 The coordinate readings have trivial common kernel
on the rational quotient. Select a finite spanning subfamily
of its dual and put the finitely many fields involved inside
one normal extension \(K'/K\). There are then \(\rho\)
independent conjugate readings \(\phi_i\) with the same
bounds as above. This selection requires no bound for
\([K':K]\).

\proofstep{The determinant of the chosen CM evaluations.}
Set 
\[
  d_\theta=\det(\phi_i(\theta_j))\ne0.
\]
For a local basis of \(Q_\lambda\) lift its elements to
\(M_\lambda\). Clearing denominators away from \(\lambda\)
preserves the finite inequality. The determinant on this
basis has order at least \(-\sum_i s_{iw}-\rho e_{w/\lambda}
b_\lambda\). Changing to the \(\theta_j\) adds precisely
\(e_{w/\lambda} L_{Q,\lambda}\) by the elementary divisor
theorem. Averaging over the places above \(\lambda\) gives
\[
 L_{Q,\lambda}\le n_\lambda+\rho b_\lambda,\qquad
 n_\lambda=
 \frac{\sum_{w\mid\lambda} f_{w/\lambda}
                    (w(d_\theta)+\sum_i s_{iw})}{[K':K]} .
\]
This signed quantity has finite support. Bounding only its
positive part would spoil the ensuing product formula.

At \(\sigma\) factor \(e^{-a_{i\sigma}+C_{\rm arch}}\) from
row \(i\) in the determinant estimate, then use the column
norm bound \(\sqrt\rho\|\theta_j\|_{\sup,\sigma|K}\).
Hadamard followed by the product formula for \(d_\theta\)
therefore gives
\[
 \langle n\rangle\le
 \rho\big(H_\theta+C_{\rm arch}+B_{\rm deg}
                              +\tfrac12\log\rho\big).
\]
When \(\rho=0\) use the empty determinant. This gives the
length on completing \(Q_\lambda\), since flatness preserves
the kernel and the finite-length quotient. In particular
both the baseline and the degree price are taken at this
lift. There is no induction on earlier local coindices.

\proofstep{The final cost in the selected rank.}
For the run under consideration \(A\le S_\kappa\), and
the selected radius and torsion estimates bound
\(C_{\rm arch}\), \(B_{\rm deg}\) and \(\langle b\rangle\)
each by \(C_0 z^4 S_\kappa\). The bounds use \(r\ge3\),
the good level, \(\log {\rm N}\mathfrak c_K\le L\) in
the size data, the positive central and last weights,
\(S_\kappa\le C_K kfb\), \(T,c>0\) and \(Tc\le1\).
For the cutoff parameter fixed at seven one may use
\(K_*=20+|C_A|+7|C_A|+|C_K|\) and

\[
  C_0=17+2K_*+200K_*^4.
\]
 The selected dimension bound
gives \(\log\rho\le1041z^3\) when \(\rho\ge1\).
With \(H_\theta\le z^4 S_\kappa\) the two estimates above
give \[
 \langle n\rangle\le\rho P_Q S_\kappa,\qquad \langle b\rangle\le P_Q S_\kappa,
\] with
\(P_Q=(1+3C_0+1041)z^4\). Their sum gives the required
\(2r_q P_Q S_\kappa\) for the fourth price. The degree uses the same line as the two coefficient bounds.
The numerical comparison for the radius and torsion order
is at the chosen weight and cutoff on this lift. \label{cor:slice-budget-2}

\end{proof}

\paragraph{Absorbing all four terms at the selected scale.}

The four bounds can now be compared at the scale fixed for both orientations. The numerical inequalities below show how their sum fits the allowed remainder.

Put \(P=C_2z^{A_2}\) with \(z=fr\log b\), and put \(P_D=H_z,\ P'_Q=2P_Q\), or increase these to the nonnegative majorants for the comparison. The numerical conditions used here are
\begingroup\small
\begin{equation}\label{eq:slice-remainder-scales}\begin{aligned}
 T S_L&\le2 S_\kappa,&
 |w_F|&\le(s+2)S_L+5fr^2T,&Tc_+g_0^*&\le1,\\
 60G_c+180P_H+44P_D+22P'_Q&\le P^2s,&
 5G_cfr^2+P_d&\le P^3,& P_\eta&\le P,\\
 S_\kappa&\le C_K kfb,&4C_K&\le P^2 f,&
             P^3 T^2&\le b,\\
 P^4s^3kf&\le Tm\delta_2,&
 1\le m&\le k,&f,T,g_0^*&\ge1 .
\end{aligned}\end{equation}
\endgroup
Here \(0<\delta_2\le1\) and \(P\ge87\). In the pre-final
case \(c_+=c\); omit the auxiliary term in the final case.
The coefficient bounds are polynomial bounds in the same
\(z\). For example the cutoff used for the field, integral
section and quotient constants includes \(z\ge
C_L+|C_K|+|C_q(C_A,C_K)|\) and \(z\ge10^9\).
It controls \(C_q z^4\) by \(z^5\). Once these fixed
constants are chosen the cutoff is independent of the run
place, component, lift, short sections, \(V\) and \(h\).
\label{cor:slice-budget-3}

\begin{established}[One order and precision for the two orientations]
\label{hyp:common-slice-scales}\label{par:slice-scales}
\hypertarget{source-slice-scales}{}
Take the selected rank of Lemma~\ref{lem:rank-window}. Write
\(L=x^\epsilon,\ b=e^L,\ z_0=frL,\ D_*=f\binom r2\).
With the exponents \(A_2=12,\ C_p=1\) of
Lemma~\ref{lem:window-parameter-choice}, put \(P=C_2z_0^{A_2}\)
for the fixed precision constant \(C_2\). On both runs, at every lift
and all sufficiently large selected scales, the order \(T=\lceil
x^{377}\rceil\) satisfies
\[
 \frac{P^4D_*^3kf\,x^{C_p}}{m\delta_2}\le T,\qquad
 P^3T^2\le\sqrt b,\qquad
 (D_*+2)n_W+2r\le T,\qquad 40(D_*+2),\,11D_*^2\le T .
\]
The cut has dimension parameter \(s=D_*\), and
\(n_P=8fr+4\le n_W\).
\label{cor:slice-scales-1}
\end{established}
\begin{proof}
For \(p+q=r\) the unitary base has dimension \(fpq\); the flags in
the two summands contribute \(f\binom p2+f\binom q2\).
Their sum is \(D_*\), including at \(q=0\).
Lemma~\ref{lem:rank-window} chooses
\[
 n_W=8\left\lceil\frac{10r}8
         \left(f+\frac{f^2(k-1)^3L^2}{m}\right)\right\rceil,
 \qquad r+t=n_P+Tn_W .
\]
In particular \(n_W\ge10fr\ge n_P\). No order is chosen anew after
the rank \(t\) has been fixed.

For an explicit comparison put \(y=x^{1/100}\). The bounds in that
lemma give \(f\ll y^{216}\), \(k,r\ll y^{470}\), \(m\ge y^{49}\),
\(L\le y^2\), and \(n_W\ll y^{2267}\). The nonempty set of larger
indices has density \(\delta_2\ge 1/n\), where \(2n=f\). Thus
\(z_0\ll y^{688}\), \(D_*\ll y^{1156}\), and
\[
 \frac{P^4D_*^3kf\,x}{m\delta_2}
 \ll_{C_2} y^{48\cdot688+3\cdot1156+470+216+100-49+216}
 =y^{37445}\ll x^{377}.
\]
The remaining lower bounds for \(T\) have still smaller exponents.
Finally \(\log(P^3T^2)=O_{C_2}(\log x)\), whereas
\(\log\sqrt b=x^\epsilon/2\). All quantities in this comparison are
fixed before the orientation and lift are chosen.
\end{proof}

\begin{lemma}[Size bounds in both orientations]
\label{lem:common-slice-scales}
For the data of Lemma~\ref{hyp:common-slice-scales}, both orientations satisfy the finite boundary estimate, the allowance conditions in \eqref{eq:slice-remainder-scales}, and the inclusion at every weight coordinate
\[
 T(w_L)_{\tau,i}\le\kappa_{\tau,i}\qquad\text{for all }\tau,i.
\]

For fixed choices of the constants in the polynomial \(P\), the
threshold on \(x\) is common to all lifts in the two runs.
\end{lemma}

\begin{proof}
\proofstep{Compatibility of the weight inclusions.}
The weight inclusion is coordinatewise:
\(T\) times the light rows is bounded by the selected
\(\kappa\) rows, so \(TD_L\le D_\kappa\).
The weight size has the differential degree, a central term
\(fr\mu_j n\), and a constant one, with \(1\le\mu_j\le r\).
The three central parameters at \(w_L,w_F,\kappa\) are
\(n_W,n_P,Tn_W+n_P\). The constant term in the weight size gives \(TS_L\le S_\kappa+T\).
Using \(T\le b\) and the head-weight bounds \(b\le S_\kappa\le7kfb\),
\[
 TS_L\le S_\kappa+T\le S_\kappa+b\le2S_\kappa .
\]

For the short weight, the row-gap
remainder after subtracting \(T w_L\) is at most
\((s+2)\) times the light gap plus \(3T\). For the
terminal gap of the large block the adjustment
\((s+2)n_W+2r\) is at most \(T\), giving that same remainder
bound. Summing the gap comparisons from the last row in
each block bounds the differential degree by
\((s+2)D_L+3fr^2T\). The short central term is at most
\(fr^2 n_P\le fr^2T\). Including the constant term gives
\(|w_F|\le(s+2)S_L+5fr^2T\).
\label{cor:slice-scales-2}

\proofstep{Room for the finite boundary term.}
For the finite boundary cost write \(Q=\min(g_0,g_1)\),
\(R_7=(\operatorname{rd}E)^7\). The record uses
\[
             c=R_7P D_L^2/Q,
             \qquad g_0^*=1250r .
\]
The level-size input is \(Q\ge(kb)^8\).
To see the large margin in this comparison, let
\(B=\sqrt b\). The rank and weight data give
\(r\le2k^2\), \(f,m\le k\) and \(D_L\le2frb\le4k^2B^2\),
where \(f\le k\) and \(r\le2k\) give the second bound.
Also \(R_7\le B\), \(P^3T^2\le B\), \(TP\le B\)
and \(B\ge10^{100}\). Thus
\[
\begin{split}
 R_7^4P^3D_L^2 &\le16 k^4 B^9,\\
 Tg_0^*R_7PD_L^2 &\le40000 k^6 B^6,\\
 2fm^3 R_7^3 P D_L^2 &\le32 k^8 B^8 .
\end{split}
\]
Each is at most \(k^8 B^{16}\le Q\).

The middle bound
gives \(Tc g_0^*\le1\), and the last pays the auxiliary
cusp allowance at the same order. The short-basis
comparison for the \hyperref[par:slice-finite-depth]{finite readings} uses
\(R_0\le R_7\) and \(N_0^4n_W+2N_0^3\le P\).
These give its two pointwise inequalities for \(c\).
Also \[
 J_0\le2\sqrt{N_0}\le3k,\qquad D_LJ_0\le12k^3 b<Q.
\] The degree estimate
\(2D_\kappa+2<g_i\), \(TD_L\le D_\kappa\) and
\(T\ge40(D_*+2)\) give the level conditions for the
Hodge--Tate comparison. The inequality \(4D_L+4\le g_i\)
also follows directly from \(D_L\le4k^2b\) and
\(g_i\ge(kb)^8\).
\label{cor:slice-scales-3}

\proofstep{Checking the four coefficient allowances.}
The coefficient bounds for the four prices use the
polynomial domination at this same \(z_0\), including
the auxiliary, height and quotient constants at their
fixed exponents. Its threshold is chosen once after
those constants have been chosen. The record supplies
\(z_0\ge x\) from \(x\le fr\) and \(L\ge1\), so a
single lower bound on \(x\) ensures the required
polynomial domination for every place and for all
lifts on either run.
In particular the single threshold \(x_*\) chosen for the comparison has the implication
\[
  x\ge x_*\quad\Longrightarrow\quad z_0\ge x_*.
\]
 We use the construction in the specified data on the common family, including the
level interval and positive density of the indices with larger weights, as a
size input here. This argument derives the two-run
inequalities used by the index on that record.
\label{cor:slice-scales-4}
\end{proof}


\begin{proposition}[The remainder at the selected scale]\label{prop:scaled-geometric-remainder}
On either selected run with \(0<\delta_2\le1\), \(P\ge87\),
the bounds \eqref{eq:slice-remainder-scales}, and the common
choice in Lemma~\ref{lem:common-slice-scales}, the weights
and branch lattices in Proposition~\ref{prop:local-product-colength}
satisfy
\[
 \langle r\rangle\le R_j fmb\,\delta_2.
\]

The choice of scale precedes the place and the tested subspace.
\end{proposition}

\begin{proof}
To check the consumption of these inequalities, use
\(R\delta\le d\) and the rank bounds in the four prices.
The light part of \(H_s\) above, after separating
\(5G_c fr^2T\), is at most
\((s+4)(s+1)(3G_c+9P_H) S_L\).
For instance the coefficient of \(G_c\) before this
simplification is
\((s+2)+2(s+4)(s+1)+(s+1)\).

For the \(P_H\) terms use \(P_{rw}\le P_H\) and
\(\tfrac43s(s+4)+\tfrac23s+3(s+1)
 \le3(s+4)(s+1)\). With \(s\ge2\),
\((s+4)(s+1)\le10s^2\). Division by \(d>0\) now gives
\begin{align*}
 \langle r\rangle/d\le{}&
 (60G_c+180P_H+44P_D+22P'_Q)s^2 S_\kappa/T\\
 &+(5G_cfr^2+P_d)T+\tfrac12 P_\eta T c_+ .
\end{align*}
For the first summand, combine the allowance
\(P^2s\) with \(S_\kappa\le C_K kfb\) and obtain
\[
 \frac{P^2s^3S_\kappa}{T}
 \le\frac{C_K}{P^2}\,mb\delta_2
 \le\tfrac14 fmb\delta_2 .
\]
The first inequality uses \(P^4s^3kf\le Tm\delta_2\);
the second uses \(4C_K\le P^2f\).
For the second summand use its coefficient bound \(P^3\)
and \(P^3T^2\le b\).

Since \(Tf m\delta_2\ge4\), the
resulting allowance \(P^3T\le b/T\) is at most
\(\tfrac14 fmb\delta_2\). Finally \(Tc_+g_0^*\le1\)
and \(P_\eta\le P\) give
\(\tfrac12P_\eta T c_+\le P^3T\).
Adding these three estimates proves the asserted remainder bound.
Rank zero has zero lengths. This yields the weighted
remainder in \eqref{eq:geometric-run-index} for either
selected run with its own degree normalization. In
particular the paired trace calculation can use its
\(12fmb\) allowance without a factor for the number of
branches, primes or choices of a subspace.
\end{proof}
